\section{Necessary root tests}
\label{sec:tests}

A root $x\in R$ is called \emph{available at time $m$} if no collision
through time $m$, excluding index $i$, has its literal on $x$ as its unique
root literal true under $w$.  A single collision excludes at most one root.
Thus all $k$ roots are available on $\{V=0\}$ and at least $k-1$ roots are
available on $\{V=1\}$.

\begin{lemma}[Robust root-test bounds]
\label{lem:root-tests}
For $i\le m<L$, conditional on $\cG_m$,
\begin{align}
    \Prb(A_m\mid\cG_m)
    &\le
      B_{\mathrm{t}}^d\ind_{\{S_m\ne\varnothing\}}p_m^k
      &&\text{on }\{V=0\},
      \label{eq:root-test-zero}\\
    \Prb(A_m\mid\cG_m)
    &\le
      B_{\mathrm{t}}^d\ind_{\{S_m\ne\varnothing\}}p_m^{k-1}
      &&\text{on }\{V=1\}.
      \label{eq:root-test-one}
\end{align}
\end{lemma}

\begin{proof}
Fix an available root $x$.  On $A_m$, the hidden nonroot parts of all test
clauses for $x$ through time $m$ must kill $S_m$.  Otherwise there is some
$z\in S_m$ satisfying every one of those test parts.  Flip only $x$ in the
assignment $(w,z)$.  Every clause of $F_m^{(-i)}$ remains satisfied:
\begin{itemize}
\item a clause with no root literal true under $w$ already has a satisfied
      nonroot literal, which is unchanged;
\item a clause with another root literal true under $w$ retains that true
      literal;
\item a clause whose unique true root literal is on $x$ cannot be a
      collision, by availability, and is therefore a test clause whose
      nonroot part is satisfied by $z$.
\end{itemize}
The flipped assignment also satisfies $C_i$, contradicting $A_m$.

Suppose root $x$ has $r_x$ test clauses through time $m$.  By the
conditional test-density hypothesis, their joint density relative to
$r_x$ independent uniform proper $(k-1)$-clauses is at most
$B_{\mathrm{t}}^{r_x}$.  Hence their probability of killing $S_m$ is at most
\[
    B_{\mathrm{t}}^{r_x}q_{r_x}(S_m)
    \le B_{\mathrm{t}}^{r_x}q_d(S_m)
    =B_{\mathrm{t}}^{r_x}p_m,
\]
because $r_x\le Z\le d$.  Test families belonging to distinct roots are
conditionally independent.  The total number of test clauses across any
selected available roots is at most $Z\le d$.  Multiplying the necessary
events for all $k$ available roots when $V=0$, or for any $k-1$ available
roots when $V=1$, proves \eqref{eq:root-test-zero} and
\eqref{eq:root-test-one}.
\end{proof}
